\documentclass[12pt,a4paper]{report}

\usepackage[utf8]{inputenc}
\usepackage[T1]{fontenc}
\usepackage[spanish]{babel}
\usepackage{geometry}
\usepackage{setspace}
\usepackage{graphicx}
\usepackage{booktabs}
\usepackage{amsmath}
\usepackage{amssymb}
\usepackage{hyperref}
\usepackage{enumitem}
\usepackage{longtable}

\geometry{
    left=3cm,
    right=2cm,
    top=3cm,
    bottom=2cm
}

\onehalfspacing

\begin{document}

% ============================================================
% PORTADA
% ============================================================

\begin{titlepage}
    \centering

    {\Large Universidad XXXXX\par}
    \vspace{0.5cm}

    {\Large Programa de Doctorado en XXXXX\par}

    \vspace{3cm}

    {\LARGE \textbf{Título provisional de la investigación}\par}

    \vspace{1cm}

    {\large
    Estimación de validez en tiempo de ejecución para la integración
    de modelos híbridos aprendidos en control de dinámica vehicular
    \par}

    \vfill

    {\large Nombre del doctorando\par}
    {\large Director / Orientador\par}

    \vspace{1cm}

    {\large Año\par}

\end{titlepage}


% ============================================================
% RESUMEN
% ============================================================

\begin{abstract}

% Escribir esta sección al final.
%
% El resumen deberá contener, en este orden:
%
% 1. Contexto:
%    - los modelos aprendidos e híbridos permiten compensar
%      discrepancias de modelos físicos en control vehicular;
%    - su utilidad closed-loop depende de que la información aprendida
%      siga siendo válida durante la operación.
%
% 2. Problema:
%    - la validez puede degradarse cuando el vehículo opera en
%      regímenes físicos, regiones de estado o condiciones de entrada
%      no representadas durante el entrenamiento;
%    - el controlador puede continuar utilizando información aprendida
%      aun cuando esta pierda confiabilidad.
%
% 3. Propuesta:
%    - modelo physics + learned residual;
%    - caracterización runtime mediante incertidumbre aleatoria,
%      incertidumbre epistémica, calibración y un indicador OOD de
%      bajo costo computacional;
%    - construcción de una señal de validez/confianza C_k;
%    - uso de C_k para adaptar un MPC/NMPC.
%
% 4. Evaluación:
%    - shifts vehiculares físicamente interpretables;
%    - comparación con baselines sin validity awareness, validity-only
%      y OOD-only;
%    - métricas de modelo, validity/OOD y, principalmente, closed-loop.
%
% 5. Contribución esperada:
%    - control-relevant validity estimation;
%    - validity-aware vehicle control;
%    - runtime validity characterization;
%    - eventualmente, taxonomía experimental de shifts vehiculares.

\end{abstract}

\tableofcontents


% ============================================================
% CAPÍTULO 1
% ============================================================

\chapter{Introducción}

% La introducción debe construir UNA única cadena causal:
%
% cambio de condición operativa
% -> cambio de validez de la información aprendida
% -> posible degradación del control
% -> estimación runtime de validez
% -> reacción del controlador.
%
% No presentar OOD + uncertainty como el ``gap'' en sí mismo.
% Son herramientas candidatas para estimar validez runtime.

\section{Contexto y motivación}

% Desarrollar desde el problema físico/controlístico:
%
% - la dinámica vehicular depende de condiciones operativas y de
%   parámetros que pueden variar durante la conducción;
% - modelos físicos control-oriented son útiles pero presentan
%   discrepancias frente a la planta real;
% - modelos híbridos physics + learning permiten aprender parte
%   del model mismatch conservando estructura física;
% - el uso de información aprendida dentro del lazo de control
%   introduce una nueva dependencia: su validez durante runtime.
%
% Evitar una introducción general a IA, vehículos autónomos o
% Machine Learning que no contribuya directamente al problema.

\section{Problema científico}

% Formular el problema central aproximadamente como:
%
% Un modelo híbrido de dinámica vehicular puede ofrecer mejores
% predicciones que un modelo físico nominal dentro del dominio
% representado por los datos. Sin embargo, cuando cambian las
% condiciones de operación, el residual aprendido puede dejar de
% representar correctamente el model mismatch. El problema no es
% únicamente detectar una muestra OOD, sino estimar si la información
% aprendida continúa siendo suficientemente válida para que un
% controlador la utilice y definir cómo debe reaccionar el control
% cuando esa validez disminuye.
%
% Definir explícitamente que la tesis estudia:
%
% (1) runtime validity characterization;
% (2) control-relevant validity estimation;
% (3) validity-aware control reaction.

\section{Delimitación del problema}

% Dominio principal:
% - dinámica vehicular;
% - modelo híbrido physics + learned residual;
% - MPC/NMPC como controlador receptor.
%
% Fuera del núcleo de la tesis, salvo conexión directa:
% - percepción;
% - detección de objetos;
% - planificación global;
% - predicción de otros agentes;
% - assurance del stack completo de autonomía.
%
% CBF: extensión opcional, no elemento obligatorio para validar
% la hipótesis principal.
%
% DAgger/reentrenamiento online: extensión opcional si el calendario
% permite estudiarlo y si la degradación observada se asocia
% principalmente a incertidumbre epistémica.

\section{Brecha de investigación}

% Esta sección se redactará después de cerrar la revisión bibliográfica.
%
% Regla de seguridad:
% NO afirmar ``nadie combina OOD y uncertainty''.
%
% La brecha debe construirse como una limitación en el cierre del ciclo:
%
% learned vehicle dynamics
% -> runtime validity
% -> control decision
% -> closed-loop behavior.
%
% Formulación provisional segura:
%
% Los modelos aprendidos son utilizados para proporcionar información
% relevante para control en aplicaciones de dinámica vehicular, y
% existen mecanismos de uncertainty estimation, OOD detection,
% robust/adaptive control y runtime monitoring. Sin embargo, la tesis
% investiga específicamente cómo caracterizar la validez runtime de un
% residual aprendido bajo shifts vehiculares físicamente interpretables
% y cómo traducir esa caracterización en modificaciones explícitas de
% un MPC/NMPC, evaluando el efecto sobre el desempeño closed-loop.
%
% Sustituir esta formulación por una versión sustentada en literatura
% una vez completada la tabla maestra del estado del arte.

\section{Pregunta de investigación}

\begin{quote}
\textbf{¿Cómo estimar en tiempo de ejecución la validez de un modelo
residual aprendido de dinámica vehicular y utilizar esa estimación para
adaptar un controlador MPC/NMPC ante condiciones de operación no
representadas durante el entrenamiento?}
\end{quote}

% Esta versión incorpora las decisiones D1--D7 sin comprometer todavía
% una técnica concreta de fusión ni una garantía formal de estabilidad.

\section{Preguntas específicas de investigación}

\begin{enumerate}[label=\textbf{P\arabic*:}]
    \item ¿Qué tipos de shifts físicamente relevantes en dinámica
    vehicular producen pérdida significativa de validez del residual
    aprendido?

    \item ¿Qué información proporcionan la incertidumbre epistémica,
    la incertidumbre aleatoria, la calibración y un indicador OOD sobre
    el error del modelo relevante para control?

    \item ¿En qué medida las señales de uncertainty, calibration y OOD
    son complementarias o redundantes para caracterizar validez runtime?

    \item ¿Cómo construir una señal de confianza/validez
    \(C_k=g(U_{\mathrm{epi}},U_{\mathrm{ale}},\mathrm{Calibration},
    \mathrm{OOD})\) que guarde relación con el error relevante para el
    controlador?

    \item ¿Cómo debe utilizarse \(C_k\) para adaptar el horizonte de un
    MPC/NMPC de manera que se reduzca la degradación closed-loop cuando
    disminuye la validez del modelo aprendido?

    \item ¿Qué mejora ofrece la arquitectura propuesta frente a un
    controlador basado únicamente en el modelo físico, un controlador
    que utiliza el modelo aprendido sin monitoreo de validez y
    arquitecturas que utilizan únicamente validity o únicamente OOD?

    \item Como extensión opcional, ¿puede la información de validez
    emplearse para modificar la acción de una CBF o activar estrategias
    de actualización de datos/modelo ante incertidumbre epistémica?
\end{enumerate}

\section{Hipótesis de investigación}

\begin{quote}
\textbf{La combinación calibrada de incertidumbre epistémica,
incertidumbre aleatoria y evidencia de distribución fuera del dominio
puede proporcionar una estimación útil de la validez runtime de un
residual aprendido de dinámica vehicular; al incorporar explícitamente
esta estimación en la adaptación de un controlador MPC/NMPC, es posible
reducir la degradación de desempeño closed-loop bajo cambios de régimen
respecto a arquitecturas que utilizan el modelo aprendido sin
validity awareness.}
\end{quote}

% Esta hipótesis debe presentarse como algo a validar, no como resultado.
% Si la revisión muestra que conviene mantener uncertainty y OOD como
% señales separadas, reformular la parte de ``combinación'' sin alterar
% el problema científico central.


% ============================================================
% CAPÍTULO 2
% ============================================================

\chapter{Objetivos}

\section{Objetivo general}

Desarrollar y evaluar una metodología de estimación de validez en tiempo
real para un modelo híbrido de dinámica vehicular basado en un modelo
físico y un residual aprendido, e integrar dicha estimación en un
controlador MPC/NMPC para reducir la degradación del desempeño
closed-loop ante condiciones de operación no representadas durante el
entrenamiento.

\section{Objetivos específicos}

\begin{enumerate}[label=\textbf{OE\arabic*:}]
    \item Formular una arquitectura híbrida de dinámica vehicular
    compuesta por un modelo físico control-oriented y un residual
    aprendido.

    \item Definir operacionalmente los tipos de shift/OOD relevantes
    para el estudio, distinguiendo al menos régimen físico no
    representado y dynamics OOD, e incluyendo input OOD si el alcance
    temporal lo permite.

    \item Estimar y calibrar incertidumbre aleatoria y epistémica del
    modelo aprendido mediante un modelo heteroscedástico y un ensemble
    pequeño, complementado con un indicador OOD de bajo costo
    computacional.

    \item Caracterizar experimentalmente la relación entre shift,
    uncertainty, calibration, OOD, error de predicción y error relevante
    para control.

    \item Definir una representación de confianza/validez runtime
    \(C_k=g(U_{\mathrm{epi}},U_{\mathrm{ale}},\mathrm{Calibration},
    \mathrm{OOD})\) y evaluar su capacidad para anticipar o representar
    degradación del modelo útil para el controlador.

    \item Integrar la señal de validez en un MPC/NMPC mediante adaptación
    de su horizonte y evaluar el efecto sobre desempeño y robustez
    empírica closed-loop.

    \item Comparar la propuesta con arquitecturas de referencia que
    separen el efecto del modelo aprendido, la estimación de validez y
    la detección OOD.

    \item Como extensión, si el cronograma lo permite, evaluar el uso de
    la señal de validez en una CBF y/o una estrategia de adquisición y
    reentrenamiento basada en DAgger para condiciones asociadas a
    incertidumbre epistémica.
\end{enumerate}


% ============================================================
% CAPÍTULO 3
% ============================================================

\chapter{Fundamentos de dinámica y control vehicular}

% Este capítulo sustituye la fundamentación genérica sobre ``sistemas
% dinámicos'' e ``IA'' del documento anterior. Debe contener solamente
% los fundamentos necesarios para comprender el problema doctoral.

\section{Dinámica vehicular relevante para control}

% Incluir únicamente lo necesario:
% - estados dinámicos seleccionados;
% - entradas de control;
% - dinámica longitudinal/lateral/yaw según el modelo elegido;
% - interacción neumático--carretera;
% - papel de \mu y del régimen de operación;
% - regiones cercanas a saturación/handling limits si se estudian.
%
% Conectar siempre con el problema de cambio de régimen y model mismatch.

\section{Modelos control-oriented y model mismatch}

% Explicar:
% - por qué se utilizan modelos físicos simplificados en control;
% - fuentes de discrepancia modelo--planta;
% - efecto de parámetros desconocidos/variables;
% - por qué un residual aprendido puede ser útil;
% - por qué ese residual puede perder validez fuera del dominio de datos.

\section{Modelos híbridos: physics + learned residual}

% Definir formalmente una arquitectura del tipo:
%
% x_{k+1} = f_{phys}(x_k,u_k) + r_\theta(x_k,u_k) + w_k
%
% donde r_theta aprende la discrepancia residual.
%
% Precisar posteriormente:
% - variables de entrada;
% - salida residual;
% - horizonte de uso;
% - frecuencia de actualización;
% - hipótesis del modelo físico.
%
% No convertir esta sección en una comparación exhaustiva de algoritmos ML.

\section{Control predictivo basado en modelo}

% Presentar MPC/NMPC solo al nivel necesario para la tesis:
% - modelo de predicción;
% - horizonte;
% - función de costo;
% - restricciones;
% - receding horizon;
% - sensibilidad a errores de modelo.
%
% Introducir el horizonte N como principal variable de adaptación
% propuesta. No fijar todavía si C bajo implica aumentar o disminuir N:
% esta relación debe justificarse en la metodología/experimentos.

\section{Control Barrier Functions como extensión opcional}

% Sección breve y explícitamente secundaria.
%
% Explicar:
% - papel de una CBF como safety filter;
% - parámetro(s) que determinan la agresividad/conservadurismo;
% - cómo C_k podría modular dichos parámetros.
%
% No prometer garantía formal de seguridad hasta demostrar que la
% arquitectura y los supuestos permiten esa afirmación.


% ============================================================
% CAPÍTULO 4
% ============================================================

\chapter{Información aprendida y validez en tiempo de ejecución}

\section{Información aprendida relevante para el controlador}

% La información aprendida principal de la tesis NO es una taxonomía
% genérica de salidas de IA. Es el residual/model mismatch incorporado
% al modelo de predicción del controlador.
%
% Situar otros roles funcionales de modelos aprendidos solo como contexto
% bibliográfico: full learned dynamics, learned tire/force models,
% estimadores de parámetros, etc.

\section{Incertidumbre predictiva}

\subsection{Incertidumbre aleatoria}

% Definición, interpretación y papel dentro del modelo heteroscedástico.
% Relacionar con ruido/variabilidad no reducible bajo los supuestos del
% modelo, evitando equipararla automáticamente con invalidity.

\subsection{Incertidumbre epistémica}

% Definición e interpretación como incertidumbre asociada al conocimiento
% del modelo/datos. Introducir el small ensemble como mecanismo candidato.
% Relacionarla con regiones poco representadas, pero no afirmar que
% epistemic uncertainty sea equivalente a OOD.

\subsection{Calibración}

% Explicar por qué no basta con producir una varianza: la incertidumbre
% debe contrastarse con la frecuencia/magnitud real de los errores.
% Definir las métricas de calibration que finalmente se utilicen.

\section{OOD y distribution shift en dinámica vehicular}

% Separar explícitamente las categorías operacionales elegidas.

\subsection{Régimen físico no representado}

% Caso principal: cambios de parámetros/condiciones físicas no presentes
% durante entrenamiento. Ejemplo prioritario: variaciones desconocidas
% de coeficiente de fricción \mu.

\subsection{Dynamics OOD}

% Regiones de estado/dinámica nunca vistas o insuficientemente cubiertas
% durante entrenamiento. Ejemplo de escenario: V_x > V_{x,\max}^{train}
% o regiones dinámicas no representadas.

\subsection{Input OOD}

% Extensión opcional si el cronograma lo permite. Definir qué entradas
% del modelo constituyen input OOD y cómo se distingue de dynamics OOD.

\subsection{Fallo o degradación sensorial}

% Escenario opcional. Incluir solo si se demuestra que produce una
% alteración relevante en la información consumida por el modelo/control.
% No mezclar automáticamente sensor fault con dynamics OOD.

\section{Validez/confianza runtime}

% Sección central del marco conceptual.
%
% Establecer explícitamente:
%
% uncertainty != OOD != validity.
%
% Proponer una señal de trabajo:
\begin{equation}
C_k = g\!\left(
U_{\mathrm{epi},k},
U_{\mathrm{ale},k},
\mathrm{Cal}_k,
S_{\mathrm{OOD},k}
\right),
\end{equation}
%
% donde C_k no se define como ``probabilidad de estar correcto'' salvo
% que posteriormente exista una calibración que justifique esa lectura.
%
% La señal debe evaluarse por su relación con error relevante para
% control, no solo por AUROC de detección.

\section{Validez relevante para control}

% Diferenciar:
% - error estadístico del predictor;
% - error dinámico acumulado/multi-step;
% - error que altera la optimización del MPC/NMPC;
% - degradación closed-loop.
%
% Esta sección es una de las claves para demostrar profundidad doctoral.


% ============================================================
% CAPÍTULO 5
% ============================================================

\chapter{Estado del arte}

\section{Estrategia de revisión de literatura}

% Documentar:
% - bases de datos;
% - período;
% - strings de búsqueda;
% - criterios de inclusión/exclusión;
% - proceso de selección;
% - uso de reviews para terminología y de primarios para evidencia;
% - tabla maestra de papers cercanos.
%
% Mantener separada la evidencia bibliográfica de las decisiones propias.

\section{Modelos aprendidos e híbridos en control de dinámica vehicular}

% Organizar por rol funcional, no por algoritmo:
% 1. full learned dynamics;
% 2. hybrid physics + learning;
% 3. residual/model-mismatch learning;
% 4. learned tire/force/parameter models.
%
% Para cada familia responder:
% ¿qué información produce el modelo y cómo modifica el controlador?

\section{Uncertainty-aware learned dynamics para control}

% Revisar trabajos con:
% - epistemic/aleatoric/predictive uncertainty;
% - calibration;
% - MPC/NMPC con learned models;
% - mecanismos robust/adaptive/stochastic relacionados.
%
% Extraer no solo técnica UQ sino cómo la incertidumbre afecta control.

\section{OOD, distribution shift y runtime monitoring}

% Revisar:
% - OOD detection for dynamics models;
% - distribution shift;
% - residual/change detection;
% - temporal shifts;
% - runtime monitors en robotics/CPS cuando el mecanismo sea transferible.
%
% Distinguir detector de shift de estimador de validez.

\section{Validity-aware y confidence-aware control}

% Buscar explícitamente trabajos donde una señal de confianza/validez:
% - cambie autoridad del learned model;
% - modifique restricciones/pesos/horizonte;
% - active fallback;
% - aumente robustificación;
% - module safety filters.
%
% Esta sección debe contener los trabajos que más amenazan/apoyan la
% contribución doctoral.

\section{Síntesis crítica y brecha}

% Construir una tabla comparativa con columnas equivalentes a:
%
% Referencia | Dominio | Learned information | Modelo AI | Controlador |
% Uncertainty | OOD | Definición de validity | Shift | Cómo validity
% modifica control | Closed-loop property | Guarantee | Validation |
% Real-time | Limitación | Aporte a la tesis.
%
% La síntesis debe responder:
%
% qué existe -> qué resuelve -> qué limitaciones deja.
%
% Solo después de esta síntesis congelar la formulación final del gap.


% ============================================================
% CAPÍTULO 6
% ============================================================

\chapter{Marco conceptual y arquitectura propuesta}

\section{Cadena shift--validity--control}

% Figura conceptual recomendada:
%
% operating-condition shift
% -> learned residual
% -> {prediction, U_epi, U_ale, calibration, OOD}
% -> C_k
% -> MPC/NMPC adaptation
% -> u_k
% -> vehicle
% -> closed-loop metrics.

\section{Modelo híbrido de predicción}

% Presentar formalmente la combinación physics + learned residual.
% La implementación concreta se detallará después, pero aquí debe quedar
% claro qué información consume el MPC/NMPC.

\section{Estimador de validez}

% Arquitectura conceptual:
% - heteroscedastic model -> aleatoric uncertainty;
% - small ensemble -> epistemic uncertainty;
% - calibration layer/assessment;
% - cheap OOD score;
% - fusión/representación C_k.
%
% La función g puede ser heurística, aprendida o basada en calibración,
% pero la selección debe justificarse por su relación con control error.

\section{Integración con MPC/NMPC}

% Decisión D6 principal:
% C_k -> adaptación del horizonte del MPC/NMPC.
%
% Definir en esta sección, cuando se cierre el diseño:
% - variable adaptada: N_k;
% - límites N_min, N_max;
% - regla N_k = h(C_k);
% - suavizado/histeresis si fuese necesario;
% - implicaciones computacionales;
% - comportamiento ante C no disponible/no confiable.
%
% No asumir a priori la dirección óptima de la adaptación.

\section{Extensiones opcionales}

\subsection{Modulación de una CBF}

% Si se implementa:
% C_k -> parámetro de agresividad/conservadurismo de la CBF.
% Especificar posteriormente qué parámetro se modifica y qué propiedad
% se pretende preservar/evaluar.

\subsection{Adquisición de datos y reentrenamiento}

% Si el tiempo lo permite y la incertidumbre dominante es epistémica:
% - utilizar eventos de baja validez para recolectar datos;
% - estudiar DAgger o estrategia equivalente;
% - evaluar offline/online según viabilidad y seguridad.
%
% Esta extensión NO es necesaria para validar la contribución principal.


% ============================================================
% CAPÍTULO 7
% ============================================================

\chapter{Metodología de investigación y validación}

\section{Estrategia general}

\begin{enumerate}[label=\textbf{E\arabic*:}]
    \item Establecer el modelo físico control-oriented y el controlador
    MPC/NMPC de referencia.

    \item Entrenar el residual aprendido y cuantificar su mejora y sus
    límites frente al modelo físico dentro del dominio nominal.

    \item Implementar estimación de incertidumbre aleatoria y epistémica,
    calibración y un score OOD de bajo costo.

    \item Diseñar shifts vehiculares físicamente interpretables y
    caracterizar cómo afectan predicción, incertidumbre, OOD y validez.

    \item Construir y validar la señal \(C_k\) respecto a error relevante
    para control.

    \item Integrar \(C_k\) en la adaptación del MPC/NMPC.

    \item Comparar baselines y propuesta bajo condiciones nominales y OOD.

    \item Evaluar robustez empírica y, si el desarrollo lo permite,
    incorporar extensiones CBF y/o DAgger.
\end{enumerate}

\section{Escenarios de shift}

\subsection{Cambio de fricción desconocida}

% Escenario prioritario D7.
% Justificar por qué variaciones de \mu representan un cambio de régimen
% físico con impacto directo sobre dinámica y residual aprendido.

\subsection{Extensión del rango de velocidad}

% Escenario prioritario D7:
% V_x > V_{x,\max}^{train}.
% Definir claramente V_{x,max}^{train} y separar extrapolación de estado
% de cambios físicos de \mu.

\subsection{Fallo/degradación de sensores}

% Escenario opcional D7.
% Definir tipo de fallo, variable afectada y por qué constituye un shift
% relevante para el modelo/control. No incluirlo si no puede aislarse
% experimentalmente.

\subsection{Input OOD}

% Escenario opcional D3.
% Añadir solo si existe tiempo suficiente y una definición operacional
% clara diferente de dynamics OOD.

\section{Modelo aprendido e incertidumbre}

% D4:
% - heteroscedastic model;
% - small ensemble;
% - aleatoric + epistemic;
% - calibration;
% - cheap OOD score.
%
% Reportar costo computacional porque el uso es runtime.

\section{Construcción y evaluación de la señal de validez}

% D5:
% C = g(U_epi, U_ale, Calibration, OOD).
%
% Evaluar al menos:
% - correlación/asociación con prediction error;
% - asociación con multi-step/model error relevante para MPC;
% - capacidad de discriminar regiones de degradación;
% - sensibilidad a shifts;
% - redundancia/complementariedad de componentes;
% - latencia/costo.

\section{Integración validity-aware con MPC/NMPC}

% D6:
% - adaptar horizonte MPC/NMPC;
% - determinar regla h(C) y límites;
% - evaluar desempeño, factibilidad y costo computacional;
% - evitar prometer recursive feasibility hasta demostrarla.

\section{Baselines y ablation study}

% D9. Comparaciones mínimas:

\begin{description}
    \item[\textbf{B0 -- Physics-only:}]
    controlador con modelo físico sin residual aprendido.

    \item[\textbf{B1 -- Learned/no-awareness:}]
    controlador con modelo physics + learned residual, sin estimación de
    validity ni detección OOD.

    \item[\textbf{B2 -- Validity/uncertainty-only:}]
    controlador con modelo aprendido y señal construida a partir de
    uncertainty/calibration, excluyendo el componente OOD explícito.

    \item[\textbf{B3 -- OOD-only:}]
    controlador con modelo aprendido y reacción basada únicamente en el
    score OOD, sin estimación probabilística de uncertainty.

    \item[\textbf{P -- Proposed:}]
    controlador con modelo physics + learned residual y señal
    \(C_k=g(U_{\mathrm{epi}},U_{\mathrm{ale}},\mathrm{Cal},OOD)\),
    utilizada para adaptar el MPC/NMPC.
\end{description}

% Si es posible, incluir una comparación ``oracle'' con conocimiento
% perfecto del shift/error solamente como upper bound experimental,
% no como baseline obligatorio de la tesis.

\section{Métricas de evaluación}

\subsection{Métricas del modelo}

% - RMSE / MAE;
% - NLL;
% - multi-step prediction error;
% - residual error.

\subsection{Métricas de uncertainty, calibration y OOD}

% - calibration error / coverage según representación elegida;
% - NLL u otras proper scoring rules;
% - AUROC / AUPR / FPR@TPR para OOD cuando corresponda;
% - correlación con error;
% - detection delay para shifts temporales;
% - costo computacional.

\subsection{Métricas closed-loop}

% Estas son las métricas principales de la tesis:
% - tracking error;
% - errores de estados dinámicos relevantes;
% - constraint violations;
% - control effort;
% - feasibility;
% - recovery after shift;
% - indicadores de estabilidad/robustez cuando sea válido;
% - computation time / sampling deadline.

\section{Criterio de robustez empírica}

% D10:
% Definir una métrica de degradación de costo/desempeño frente a shift.
% La idea a evaluar será:
\begin{equation}
\Delta J_{\mathrm{prop}} < \Delta J_{\mathrm{baseline}},
\end{equation}
% donde Delta J mide el aumento de costo o degradación respecto a la
% condición nominal. Definir J y Delta J antes de los experimentos.
%
% Esta relación debe tratarse como criterio empírico, no como proof de
% robustez matemática salvo desarrollo teórico adicional.

\section{Nivel de garantía y alcance de las conclusiones}

% D10:
% Compromiso principal:
% - evaluación experimental rigurosa;
% - robustez empírica.
%
% Posible extensión:
% - probabilistic calibration/bounds, si da tiempo y los supuestos son
%   defendibles.
%
% No prometer todavía:
% - estabilidad formal;
% - recursive feasibility;
% - constraint satisfaction garantizada;
% - safety formal via CBF.
%
% Si alguna de estas propiedades se demuestra durante el doctorado,
% actualizar el alcance en la versión final de la tesis.

\section{Análisis estadístico y reproducibilidad}

% Definir antes de ejecutar la campaña final:
% - número de repeticiones/semillas;
% - intervalos de confianza;
% - comparación pareada entre arquitecturas;
% - reporte de fallos/violaciones, no solo promedios;
% - criterios de éxito/fallo;
% - costo computacional y deadline de ejecución.


% ============================================================
% CAPÍTULO 8
% ============================================================

\chapter{Resultados preliminares}

% Incluir únicamente evidencia ya obtenida y directamente conectada con
% la viabilidad de la tesis.
%
% Una estructura útil sería:

\section{Viabilidad del modelado aprendido de dinámica vehicular}

% Resultados que muestren que el residual/modelo aprendido puede capturar
% parte del model mismatch en condiciones nominales.

\section{Comportamiento de uncertainty y OOD bajo shifts}

% Solo si ya existen resultados suficientemente maduros.
% Presentar casos donde las señales responden o fallan, no solo resultados
% favorables. Relacionar con D3--D5.

\section{Integración preliminar con control}

% Si existe una implementación cerrada o parcial, mostrar:
% - arquitectura;
% - métricas closed-loop;
% - costo computacional;
% - limitaciones actuales.

\section{Discusión de resultados preliminares}

% Separar explícitamente:
% - evidencia ya demostrada;
% - indicios;
% - hipótesis todavía no verificadas;
% - decisiones metodológicas que los resultados obligan a revisar.


% ============================================================
% CAPÍTULO 9
% ============================================================

\chapter{Contribuciones esperadas}

% D8. Presentarlas como contribuciones candidatas/esperadas hasta que la
% evidencia bibliográfica y experimental permita congelarlas.

\section{Control-relevant validity estimation}

% Método/criterio para estimar si el residual aprendido conserva validez
% suficiente para su uso por el controlador, utilizando uncertainty,
% calibration y OOD con relación explícita al error relevante para control.

\section{Validity-aware vehicle control}

% Arquitectura que traduzca la estimación de validez en una modificación
% explícita del MPC/NMPC y demuestre su efecto closed-loop.

\section{Runtime validity characterization}

% Caracterización sistemática de cómo diferentes shifts vehiculares
% afectan residual error, uncertainty, calibration, OOD y comportamiento
% closed-loop.

\section{Taxonomía experimental de shifts vehiculares}

% Contribución opcional si la campaña experimental alcanza suficiente
% amplitud y profundidad. No prometerla como contribución principal en la
% qualificação si todavía no existe evidencia suficiente.

\section{Extensiones potenciales}

% - modulación validity-aware de CBF;
% - DAgger/reentrenamiento orientado por incertidumbre epistémica;
% - probabilistic bounds/calibration más fuertes.
%
% Presentarlas como extensiones, no como requisitos para que la tesis sea
% científicamente completa.


% ============================================================
% CAPÍTULO 10
% ============================================================

\chapter{Plan de trabajo}

\section{Actividades restantes}

% Relacionar actividades con OE1--OE8.
%
% Sugerencia de bloques:
% A1. cerrar revisión y gap;
% A2. formalizar modelo physics + residual;
% A3. consolidar UQ + calibration + OOD;
% A4. construir C_k;
% A5. integrar MPC/NMPC adaptive horizon;
% A6. campaña de shifts y baselines;
% A7. robustez empírica + análisis estadístico;
% A8. extensiones opcionales;
% A9. publicaciones;
% A10. redacción de tesis.

\section{Cronograma}

% Insertar cronograma realista del tiempo restante.
% Marcar CBF, DAgger, input OOD, sensor fault y probabilistic bounds como
% actividades opcionales o condicionadas al cumplimiento de hitos
% principales.

\section{Hitos de decisión}

% H1. gap y pregunta final sustentados bibliográficamente.
% H2. residual + uncertainty calibrada operando en nominal/OOD.
% H3. C_k validada respecto a control-relevant error.
% H4. MPC/NMPC validity-aware comparado contra B0--B3.
% H5. campaña final y robustez empírica.
% H6. decisión go/no-go para extensiones CBF/DAgger/bounds.


% ============================================================
% CAPÍTULO 11
% ============================================================

\chapter{Riesgos, limitaciones y estrategias de mitigación}

\section{Riesgo de redundancia entre uncertainty y OOD}

% Mitigación:
% realizar ablaciones B2/B3 y estudiar complementariedad antes de fijar g.
% Si son redundantes, la contribución puede mantenerse como validity-aware
% control con la representación mínima que resulte informativa.

\section{Riesgo de mala calibración o mala correlación con error}

% Mitigación:
% evaluar calibración explícitamente y no asumir que variance = validity.
% Considerar recalibración o representación alternativa de C.

\section{Riesgo de que la adaptación del horizonte no mejore closed-loop}

% Mitigación:
% formularlo como hipótesis experimental; comparar reglas h(C), límites y
% condiciones de aplicación. La contribución debe incluir resultados
% negativos o condiciones de validez si son científicamente relevantes.

\section{Riesgo computacional}

% Mitigación:
% small ensemble + cheap OOD score; medir inference/optimization time y
% sampling deadline. Reducir complejidad si runtime no es viable.

\section{Riesgo de sobreextensión del alcance}

% Mitigación:
% núcleo obligatorio: physics+residual + UQ/calibration/OOD + C + MPC/NMPC
% + baselines + shifts principales.
% Opcional: CBF, DAgger, input OOD, sensor fault, probabilistic bounds.

\section{Limitaciones de garantías formales}

% La tesis se compromete inicialmente con evaluación experimental rigurosa
% y robustez empírica. Explicitar que esto no equivale a estabilidad o
% seguridad formal.


% ============================================================
% CAPÍTULO 12
% ============================================================

\chapter{Conclusiones del examen de calificación}

% No son conclusiones finales de la tesis.
%
% Resumir:
% 1. problema: pérdida de validez runtime de residual aprendido;
% 2. gap sustentado por la revisión;
% 3. pregunta e hipótesis;
% 4. arquitectura physics+residual + validity estimator + MPC/NMPC;
% 5. shifts prioritarios;
% 6. estrategia experimental y baselines;
% 7. evidencia preliminar disponible;
% 8. contribuciones esperadas;
% 9. límites del alcance y extensiones opcionales.


% ============================================================
% REFERENCIAS
% ============================================================

\bibliographystyle{plain}
\bibliography{referencias}

\end{document}
